\documentclass[12pt,twoside]{report}

%%%%%%%%%%%%%%%%%%%%%%%%%%%%%%%%%%%%%%%%%%%%%%%%%%%%%%%%%%%%%%%%%%%%%%%%%%%%%

% Definitions for the title page
% Edit these to provide the correct information
% e.g. \newcommand{\reportauthor}{Timothy Kimber}

\newcommand{\reporttitle}{Analysing property sales data using Data Science}
\newcommand{\reportauthor}{Wenxiang Luo}
\newcommand{\supervisor}{Chiraag Lala}
\newcommand{\degreetype}{MSc Computing}

%%%%%%%%%%%%%%%%%%%%%%%%%%%%%%%%%%%%%%%%%%%%%%%%%%%%%%%%%%%%%%%%%%%%%%%%%%%%%

% load some definitions and default packages
%%%%%%%%%%%%%%%%%%%%%%%%%%%%%%%%%%%%%%%%%
% University Assignment Title Page 
% LaTeX Template
% Version 1.0 (27/12/12)
%
% This template has been downloaded from:
% http://www.LaTeXTemplates.com
%
% Original author:
% WikiBooks (http://en.wikibooks.org/wiki/LaTeX/Title_Creation)
%
% License:
% CC BY-NC-SA 3.0 (http://creativecommons.org/licenses/by-nc-sa/3.0/)
% 
%
%%%%%%%%%%%%%%%%%%%%%%%%%%%%%%%%%%%%%%%%%
%----------------------------------------------------------------------------------------
%	PACKAGES AND OTHER DOCUMENT CONFIGURATIONS
%----------------------------------------------------------------------------------------
\usepackage[a4paper,hmargin=2.8cm,vmargin=2.0cm,includeheadfoot]{geometry}
\usepackage{textpos}
\usepackage[numbers]{natbib} % for bibliography
\usepackage{tabularx,longtable,multirow,subfigure,caption}%hangcaption
\usepackage{fncylab} %formatting of labels
\usepackage{fancyhdr} % page layout
\usepackage{url} % URLs
\usepackage[english]{babel}
\usepackage{amsmath}
\usepackage{graphicx}
\usepackage{dsfont}
\usepackage{epstopdf} % automatically replace .eps with .pdf in graphics
\usepackage{backref} % needed for citations
\usepackage{array}
\usepackage{latexsym}
\usepackage[pdftex,pagebackref,hypertexnames=false,colorlinks]{hyperref} % provide links in pdf

\hypersetup{pdftitle={},
  pdfsubject={}, 
  pdfauthor={},
  pdfkeywords={}, 
  pdfstartview=FitH,
  pdfpagemode={UseOutlines},% None, FullScreen, UseOutlines
  bookmarksnumbered=true, bookmarksopen=true, colorlinks,
    citecolor=black,%
    filecolor=black,%
    linkcolor=blue,%
    urlcolor=black}

\usepackage[all]{hypcap}


%\usepackage{color}
%\usepackage[tight,ugly]{units}
%\usepackage{float}
%\usepackage{tcolorbox}
%\usepackage[colorinlistoftodos]{todonotes}
% \usepackage{ntheorem}
% \theoremstyle{break}
% \newtheorem{lemma}{Lemma}
% \newtheorem{theorem}{Theorem}
% \newtheorem{remark}{Remark}
% \newtheorem{definition}{Definition}
% \newtheorem{proof}{Proof}


%%% Default fonts
\renewcommand*{\rmdefault}{bch}
\renewcommand*{\ttdefault}{cmtt}



%%% Default settings (page layout)
\setlength{\parindent}{0em}  % indentation of paragraph

\setlength{\headheight}{14.5pt}
\pagestyle{fancy}
\renewcommand{\chaptermark}[1]{\markboth{\chaptername\ \thechapter.\ #1}{}} 

\fancyfoot[ER,OL]{\sffamily\textbf{\thepage}}%Page no. in the left on odd pages and on right on even pages
\fancyfoot[OC,EC]{\sffamily }
\renewcommand{\headrulewidth}{0.1pt}
\renewcommand{\footrulewidth}{0.1pt}
\captionsetup{margin=10pt,font=small,labelfont=bf}


%--- chapter heading

\def\@makechapterhead#1{%
  \vspace*{10\p@}%
  {\parindent \z@ \raggedright \sffamily
    \interlinepenalty\@M
    \Huge\bfseries \thechapter \space\space #1\par\nobreak
    \vskip 30\p@
  }}

%---chapter heading for \chapter*  
\def\@makeschapterhead#1{%
  \vspace*{10\p@}%
  {\parindent \z@ \raggedright
    \sffamily
    \interlinepenalty\@M
    \Huge \bfseries  #1\par\nobreak
    \vskip 30\p@
  }}

\allowdisplaybreaks

% load some macros
% Here, you can define your own macros. Some examples are given below.

\newcommand{\R}[0]{\mathds{R}} % real numbers
\newcommand{\Z}[0]{\mathds{Z}} % integers
\newcommand{\N}[0]{\mathds{N}} % natural numbers
\newcommand{\C}[0]{\mathds{C}} % complex numbers
\renewcommand{\vec}[1]{{\boldsymbol{{#1}}}} % vector
\newcommand{\mat}[1]{{\boldsymbol{{#1}}}} % matrix


\date{June 2022}

\begin{document}

% load title page
% Last modification: 2015-08-17 (Marc Deisenroth)
\begin{titlepage}

\newcommand{\HRule}{\rule{\linewidth}{0.5mm}} % Defines a new command for the horizontal lines, change thickness here


%----------------------------------------------------------------------------------------
%	LOGO SECTION
%----------------------------------------------------------------------------------------

\includegraphics[width = 4cm]{./figures/imperial}\\[0.5cm] 

\center % Center remainder of the page

%----------------------------------------------------------------------------------------
%	HEADING SECTIONS
%----------------------------------------------------------------------------------------

\textsc{\Large Imperial College London}\\[0.5cm] 
\textsc{\large Department of Computing}\\[0.5cm] 

%----------------------------------------------------------------------------------------
%	TITLE SECTION
%----------------------------------------------------------------------------------------

\HRule \\[0.4cm]
{ \huge \bfseries \reporttitle}\\ % Title of your document
\HRule \\[1.5cm]
 
%----------------------------------------------------------------------------------------
%	AUTHOR SECTION
%----------------------------------------------------------------------------------------

\begin{minipage}{0.4\textwidth}
\begin{flushleft} \large
\emph{Author:}\\
\reportauthor % Your name
\end{flushleft}
\end{minipage}
~
\begin{minipage}{0.4\textwidth}
\begin{flushright} \large
\emph{Supervisor:} \\
\supervisor % Supervisor's Name
\end{flushright}
\end{minipage}\\[4cm]


%----------------------------------------------------------------------------------------
%	FOOTER & DATE SECTION
%----------------------------------------------------------------------------------------
\vfill % Fill the rest of the page with whitespace
Submitted in partial fulfillment of the requirements for the MSc degree in
\degreetype~of Imperial College London\\[0.5cm]

\makeatletter
\@date 
\makeatother


\end{titlepage}



% page numbering etc.
\pagenumbering{roman}
\clearpage{\pagestyle{empty}\cleardoublepage}
\setcounter{page}{1}
\pagestyle{plain}

%%%%%%%%%%%%%%%%%%%%%%%%%%%%%%%%%%%%
%\begin{abstract}
%Your abstract.
%\end{abstract}

\cleardoublepage
%%%%%%%%%%%%%%%%%%%%%%%%%%%%%%%%%%%%
%\section*{Acknowledgments}
%Comment this out if not needed.

\clearpage{\pagestyle{empty}\cleardoublepage}

%%%%%%%%%%%%%%%%%%%%%%%%%%%%%%%%%%%%
%--- table of contents
\fancyhead[RE,LO]{\sffamily {Table of Contents}}
\tableofcontents 


\clearpage{\pagestyle{empty}\cleardoublepage}
\pagenumbering{arabic}
\setcounter{page}{1}
\fancyhead[LE,RO]{\slshape \rightmark}
\fancyhead[LO,RE]{\slshape \leftmark}

%%%%%%%%%%%%%%%%%%%%%%%%%%%%%%%%%%%%
\chapter{Introduction}

Nowadays, there are a significant amount of data generated every second. The daily lives of humans are producing it, and some other fields are generating numerous data, such as research, health care, commercial activities, and environmental information from all kinds of sensors. With these considerable amounts of data,  acquiring the relationship between some aspects of records or obtaining the patterns behind these data would benefit the whole world. For instance, new causes of diseases might be discovered, and technology development could be accelerated.
\\

However, this extensive data can be one of the main obstacles for analysis as it is approximately impossible for humans to obtain insights into the data manually. Under this circumstance, artificial intelligence (AI), a technique that empowers the computer to imitate human intelligence and manner,  could be one of the methods to mitigate this problem. It can extract patterns from large datasets and use them to make predictions based on future data and even distinguish which parts of data cause the results.
\\

In this project, some AI techniques are used on property and demographic data to gain insights and understand the factors that impact a homeowner's propensity to sell. The factors could be the distance to schools, hospitals, or supermarkets, the accessibility to public transport, and the types of the properties (flats or houses). It could be highly advantageous to estate agents who would discover homeowners with more potential to become clients and provide them with business. 

%%%%%%%%%%%%%%%%%%%%%%%%%%%%%%%%%%%%
\chapter{Literal Review}
Python is one of the most popular programming languages in the world as it is easy to develop and there are extensive packages for various functionalities. In this project, Python and its several packages would be used for loading data, preprocessing data, building and evaluating machine learning models. 

\section{Machine Learning}
Machine Learning (ML), a subset of AI, is a technique that the computer can learn and improve from data without explicit programming. The reason for using ML is that its performance is better than the traditional approach sometimes. For example, ML techniques would simplify the solution to a problem that comprises a long list of rules (spam mail detection).
\\

Supervised learning is one of the ML categories, and a model of this type is implemented for analyzing data. 

\section{NumPy}
Numerical Python (\textit{\textbf{NumPy}}) is a scientific computing package utilizing an optimized C/C++ API to use less computation time than pure Python calculations \citep{RN6}. It provides some data structures, one of the most important structures is \textbf{\emph{ndarray}}. Unlike Python lists, NumPy arrays are homogenous, meaning all the elements in them are of the same type, and their sizes are fixed \citep{RN4}.  Moreover, NumPy provides various built-in functions for different purposes, such as statistics, linear algebra,  transforms, and element-wise operation. 
\\

In this project, \textit{\textbf{Numpy}} is used as part of data preprocessing. For example, converting categorical data into integers as ML models can only use numerical as inputs or standardizing features so that they are in the range of 0 and 1. 

\section{Pandas}
The \textbf{\textit{pandas}} is a fast and compelling package capable of handling data of various types (numerical values, strings, and time) and from multiple sources (CSV, Excel, and MqSQL database). One of the \textbf{\textit{pandas}} data structures is \textbf{\textit{DataFrame}} which is appropriate to handle tabular data whose columns are of different types. Also, it could manage various operations, such as manipulating missing values, creating pivot tables (summary of the table), and grouping data from different columns \citep{RN4}. 
\\

In this project, \textbf{\textit{pandas}} is used for data loading and preprocessing. First, data is loaded, and it is represented as \textbf{\textit{DataFrame}}. Then the attributes of the data, such as distribution, should be inspected. Finally, the data is preprocessed, and one of the steps is handling missing values. For example, they can be dropped or filled with the mean values. 

\section{Scikit-learn}

\section{PyTorch}
\textbf{\textit{PyTorch}} was developed by Facebook, and it is one of the popular libraries for building ML models. It is also the basis of numerous packages that the developers could take advantage of \citep{RN5}. Moreover, this library could perform auto differentiation by using graphic processing unit (GPU) acceleration, which results in less training and evaluation time. 
\\

In this project, some of the components provided by \textbf{\textit{PyTorch}} are utilized.
\begin{itemize}
	\item \textbf{\textit{nn.Sequential}} is the container that contains all the layers of the ML model
	\item An appropriate loss function would be selected to compute the loss and perform backpropagation. 
	\item One of the optimizers would be used to update the model's parameters. 
\end{itemize}

\section{Matplotlib}
\textbf{\textit{Matplotlib}} is one of the Python plotting packages that can plot different types of figures, such as histograms, pie plots, and line plots. In this project, this library is primarily used for data visualization. For instance, the distribution of the raw data should be visualized to determine the procedures of preprocessing and plotting the training loss and validation loss during model evaluation. 



%%%%%%%%%%%%%%%%%%%%%%%%%%%%%%%%%%%%
\chapter{Project Plan}

\section{Data Preprocessing}

\section{Build Model with PyTorch}
\subsection{Model architecture}
\subsection{Hyperparameters Tuning}

\section{Evaluation}


%%%%%%%%%%%%%%%%%%%%%%%%%%%%%%%%%%%%
%\chapter{Experimental Results}


%%%%%%%%%%%%%%%%%%%%%%%%%%%%%%%%%%%%
%\chapter{Conclusion}


%% bibliography
\bibliographystyle{apalike}
\bibliography{proposal}

\end{document}
